\begin{frame}[fragile]{Duplicate SSA Uses}
  These two expressions are \textbf{algebraically equivalent}, and in both cases \texttt{y} holds an \emph{even} number. Does the compiler have the right to rewrite the expression?
    \begin{columns}[c]
        \begin{column}{0.44\textwidth}
            \begin{lstlisting}[style=llvmmono]
; y = 2x                
%y = mul i32 2, %x
            \end{lstlisting}
        \end{column}
        \begin{column}{0.08\textwidth}
            \only<1>{\transform}
           \only<2>{\illegaltransform} 
        \end{column}
        \begin{column}{0.44\textwidth}
            \begin{lstlisting}[style=llvmmono]
; y = x + x
%y = add i32 %x, %x
            \end{lstlisting}
        \end{column}
    \end{columns}
    \only<2>{
      \vspace{1.5em}
      \textbf{No}: if \texttt{x} is \textbf{undef} then in the second case \texttt{y} could be any legal \texttt{i32} value.\\
    However, LLVM \textbf{incorrectly} performs similar transformations
    }
    \pdfpcnote{
**First Reveal**<br/>
The two expressions are algebraically equivalent over integers. Ask the audience:
does the compiler have the right to perform this rewrite?<br/>
**Second Reveal**<br/>
No. If x is undef, mul 2 undef picks one arbitrary value and doubles it — always
even. But add undef undef reads undef twice, each time getting a potentially
different value, so the result can be odd. LLVM incorrectly performed this
rewrite. This is one of the key motivations for removing undef entirely.
}
\end{frame}
